In this chapter, we introduce the problem studied in this dissertation as an application of the Calculus of Variations and derive the corresponding mathematical model. We then apply the Direct Method of the Calculus of Variations to establish the existence of a minimiser and derive the associated Euler--Lagrange equation. Throughout this chapter, the problem is formulated using the arc-length parameter $s$, leading to the variational problem
\begin{equation}
\label{eq:ModelPrblem_s}
\tag{P}
J[\gamma(s)]
=
\underset{\gamma \in \Omega}{\min} \int_{s_0}^{s_1}
\left(
c(\gamma(s))
+
\lambda |\gamma''(s)|^2
+
\mu
\right) ds.
\end{equation}

\section{Mathematical modelling}
\label{sec:problem_formulation}

We consider a path-planning problem in a two-dimensional domain
representing a terrain. The objective is to determine a path between two
boundary points while balancing three competing considerations: the
cost associated with the terrain traversed, the cost from the length of the
path, and the cost of a curvature penality along the path.

Let $\Omega\subset\mathbb{R}^2$ denote the domain representing the
terrain, and let $a,b\in\Omega$ denote the prescribed starting and
ending points, respectively. An admissible path connecting these
points is represented by a sufficiently smooth curve
\[
\gamma:[s_0,s_1]\rightarrow\Omega,
\qquad
\gamma(s)=(x(s),y(s)),
\]
satisfying the boundary conditions
\[
\gamma(s_0)=a,
\qquad
\gamma(s_1)=b.
\]
Here, $s$ denotes the arc-length parameter along the curve.

The collection of admissible paths satisfying these boundary
conditions, together with any additional constraints imposed by the
terrain, is denoted by $\Gamma$.

The terrain is characterised by a cost field

$$
c:\Omega\rightarrow\mathbb{R}, \qquad \Omega \subset \mathbb{R}^2.
$$

where $c(x,y)$ represents the cost associated with traversing the
terrain at the point $(x,y)$. Regions with larger values of $c(x,y)$
are therefore more costly for the path to traverse. The cost incurred
by a path due solely to the terrain is given by the line integral
\begin{equation}
J_{\mathrm{terrain}}(\gamma) = 
\int_\gamma c(x,y),ds.
\label{eq:terrain_functional}
\end{equation}

While minimising this quantity favours paths through regions of low
terrain cost, it does not by itself account for the cost induced by length or curvature
of the path. A path could, for example, reduce its terrain
cost by taking a substantially longer route. Similarly, a path with
large changes in direction may be undesirable even if it has low
terrain cost and relatively short length as this is not realisic for a railway. We incorporate
additional penalties for path length and curvature.

The resulting functional is defined by
\begin{equation}
\boxed{
J(\gamma) = 
\int_\gamma
\left[
c(x,y)
+
\lambda\kappa^2
+
\mu
\right] ds,
}
\label{eq:ModelProblem}
\end{equation}
where $\lambda$, and $\mu$ are non-negative parameters controlling
the contribution of the terms of the integrand. These terms represent terrain cost, curvature, and length,
respectively.

The path-planning problem is consequently formulated as the
minimisation problem
\begin{equation}
\inf_{\gamma\in\Gamma}J(\gamma).
\label{eq:extended_problem}
\end{equation}

We now reformulate the integrand in \eqref{eq:ModelProblem} term by term to obtain the more convenient formulation given in \eqref{eq:ModelPrblem_s}.


\subsection{Terrain Cost Term}
\label{subsec:Terrain_term}

The first term in the integrand of \eqref{eq:ModelProblem} represents
the cost associated with traversing the terrain. Let
$$
\gamma(s)=(x(s),y(s)),
$$
denote the path parametrised by arc length $s$. Since the cost field is
defined by
$$
c:\Omega\rightarrow\mathbb{R},
$$
the cost at any point on the path is obtained by evaluating $c$ at the
position of the curve. Thus,
$$
c(x,y)
\quad\longrightarrow\quad
c(x(s),y(s))
=
c(\gamma(s)).
$$

The key step here is the composition of the cost field
with the path,
$$
c\circ\gamma:[s_0,s_1]\rightarrow\mathbb{R},
\qquad
(c\circ\gamma)(s)=c(\gamma(s)).
$$

In other words, the terrain cost along the path is represented by the
composition $c\circ\gamma$. This makes explicit why the cost field
$c(x,y)$ becomes $c(\gamma(s))$ when the path is parametrised by $s$.
rather than evaluating the cost at an arbitrary point in the domain,
we evaluate it at the position of the curve as it traverses the
terrain.

Consequently, the terrain component of the functional becomes

$$
\int_\gamma c(x,y)\,ds
=
\int_{s_0}^{s_1}c(\gamma(s))\,ds.
$$

\subsection{Path Length Term}
\label{subsec:length_penalty}

The final term in the integrand of \eqref{eq:ModelProblem} represents a
penalty on the total length of the path. For a curve

$$
\gamma:[s_0,s_1]\rightarrow\mathbb{R}^2,
$$
the length of the path is given by
\begin{equation}
L(\gamma)
\int_{s_1}^{s_0}\gamma \, ds,
\label{eq:path_length}
\end{equation}
where $s$ denotes the arc-length parameter.

More generally, if the curve is initially parametrised by an arbitrary parameter  
$t\in[t_0,t_1]$, its length is given by
$$
L(\gamma)
=
\int_{t_0}^{t_1}
|\gamma'(t)|\,dt.
$$
Introducing the arc-length parameter
$$
s(t)
=
\int_{t_0}^{t}
|\gamma'(\tau)|\,d\tau,
$$
allows the curve to be parametrised directly by the distance travelled
along the path. Consequently,
$$
|\frac{d\gamma}{ds}|=1.
$$

Thus, for an arc-length parametrised curve, an change in $ds$
corresponds directly to an change in the length of the path, giving
$$
L(\gamma)
=
\int_{s_0}^{s_1}1\,ds.
$$
The length penalty in the functional is therefore
\begin{equation}
J_{\mathrm{length}}(\gamma) =   \mu L(\gamma) =
\mu\int_{s_0}^{s_1}1,ds.
\label{eq:length_functional}
\end{equation}
Hence, the contribution of the length penalty to the integrand of
\eqref{eq:ModelProblem} is simply
$$
\mu.
$$
The parameter $\mu\geq0$ controls the relative importance assigned to
the total length of the path. Increasing $\mu$ therefore places a
greater penalty on longer paths.


\subsection{Curvature Cost Term}
\label{subsec:curvature_penalty}

The second term in the integrand of \eqref{eq:ModelProblem} penalises
the curvature of the path. To express this term in terms of the
derivatives of $\gamma$, we first consider a general parametrisation

$$
\gamma(t)=(x(t),y(t)),
\qquad t\in[t_0,t_1].
$$
For a sufficiently smooth planar curve, the curvature is given by
\begin{equation}
\kappa(t) = 
\frac{\left|\gamma'(t)\times\gamma''(t)\right|}
{\left|\gamma'(t)\right|^3}.
\label{eq:curvature_general}
\end{equation}

We now consider the arc-length parametrisation used in the formulation
of the functional. As shown in the previous subsection, when $s$ is
the arc-length parameter, the curve has unit speed, so that
\begin{equation}
\left|\gamma'(s)\right|=1.
\label{eq:unit_speed}
\end{equation}
It can be shown that, $\gamma'(s)$ and $\gamma''(s)$ are orthogonal. Consequently,

$$
\left|\gamma'(s)\times\gamma''(s)\right|
=
\left|\gamma'(s)\right|
\left|\gamma''(s)\right|
=
\left|\gamma''(s)\right|,
$$
where \eqref{eq:unit_speed} has been used. Substituting these results into \eqref{eq:curvature_general} gives
$$
\kappa(s)
=
\frac{\left|\gamma''(s)\right|}
{\left|\gamma'(s)\right|^3}
=\frac{\left|\gamma''(s)\right|}
{(1)^3} = 
\left|\gamma''(s)\right|.
$$
Therefore,
\begin{equation}
\kappa(s)=\left|\gamma''(s)\right|.
\label{eq:curvature_arclength}
\end{equation}
The curvature contribution to the functional is consequently
\begin{equation}
J_{\mathrm{curvature}}(\gamma)=
\lambda\int_\gamma\kappa^2 ds.
\label{eq:curvature_functional}
\end{equation}
Using \eqref{eq:curvature_arclength}, this becomes
$$
J_{\mathrm{curvature}}(\gamma)
=
\lambda
\int_{s_0}^{s_1}
\left|\gamma''(s)\right|^2 ds.
$$
Hence, the curvature contribution to the integrand of
\eqref{eq:ModelProblem} can be written as

$$
\lambda\kappa^2
=
\lambda\left|\gamma''(s)\right|^2.
$$
The use of $\kappa^2$ penalizes curvature independently of its
direction, while assigning a greater cost to regions where the path
bends more sharply. The parameter $\lambda\geq0$ therefore controls the
relative importance of path smoothness within the functional.

\section{Euler--Lagrange Equations - Classical Approach}

As established in the previous section, the proposed path-planning problem can be expressed in terms of the arc-length parameter $s$. Recalling \eqref{eq:ModelPrblem_s}, the functional is given by
\begin{equation}
J[\gamma]
=
\int_{s_0}^{s_1}
\left(
c(\gamma(s))
+
\lambda |\gamma''(s)|^2
+
\mu
\right) ds.
\end{equation}

The corresponding Euler--Lagrange equation provides a necessary condition for a sufficiently smooth minimiser of the functional. For a second-order functional of the form
$$
J[\gamma]
=
\int_{s_0}^{s_1}
F(\gamma,\gamma',\gamma'')\,ds,
$$
the second-order Euler--Lagrange equation is given by
\begin{equation}
\frac{\partial F}{\partial\gamma}
-
\frac{d}{ds}
\left(
\frac{\partial F}{\partial\gamma'}
\right)
+
\frac{d^2}{ds^2}
\left(
\frac{\partial F}{\partial\gamma''}
\right)
=0.
\end{equation}
For the functional $F(\gamma,\gamma',\gamma'')= c(\gamma)
+\lambda|\gamma''|^2
+\mu$, we obtain
$$
\frac{\partial F}{\partial\gamma}
=
\nabla c(\gamma),
\qquad
\frac{\partial F}{\partial\gamma'}
=
0,
\qquad
\frac{\partial F}{\partial\gamma''}
=
2\lambda\gamma''.
$$
Substitution into the Euler--Lagrange equation therefore gives
$$
\nabla c(\gamma)
+
2\lambda\gamma^{(4)}
=
0,
$$
or equivalently
\begin{equation}
\label{eq:ModelProblem_EL_Solution}
\gamma^{(4)}(s)
=
-\frac{1}{2\lambda}
\nabla c(\gamma(s)).
\end{equation}
Equation \eqref{eq:ModelProblem_EL_Solution} characterises the classical stationary paths associated with the proposed functional. In general, however, solving this equation explicitly is not straightforward. In particular, this equation is a fourth-order nonlinear system of ordinary differential equations.




\subsection{Gaussian Hill}
To illustrate the resulting nonlinearity, consider a uniform background cost containing a region of increased cost. A Gaussian hill provides a simple model for such a region, given by
\begin{equation}
\label{eq:hill_C}
c(\gamma(s))
=
H\exp\left(
-\frac{|\gamma(s)-p|^2}{2\sigma^2}
\right),
\end{equation}
where $H>0$ represents the magnitude of the increase in cost, $p\in\mathbb{R}^2$ is the centre of the hill and $\sigma>0$ controls its spatial spread. If the background cost is $c_0$ and the maximum cost is $c_{\max}$, then
$$
H=c_{\max}-c_0.
$$

The gradient of \eqref{eq:hill_C} with respect to $\gamma$ is
$$
\nabla c(\gamma)
=
-\frac{H}{\sigma^2}
\exp\left(
-\frac{|\gamma-p|^2}{2\sigma^2}
\right)
(\gamma-p).
$$
Therefore, equation \eqref{eq:ModelProblem_EL_Solution} becomes
\begin{equation}
\label{eq:E_L_Gaussian}
\gamma^{(4)}(s)
=
\frac{H}{2\lambda\sigma^2}
\exp\left(
-\frac{|\gamma(s)-p|^2}{2\sigma^2}
\right)
(\gamma(s)-p).
\end{equation}
Even for this relatively simple cost landscape, the equation \eqref{eq:E_L_Gaussian} is a nonlinear fourth-order system. Consequently, obtaining an explicit analytical solution is generally difficult and would require solving a nonlinear boundary-value problem. Rather than attempting to obtain a closed-form solution, the remainder of this chapter therefore investigates the problem from a variational perspective. In particular, we shall use the Direct Method of the Calculus of Variations to establish conditions which the functional admits a minimiser, without requiring that minimiser to be explicitly constructed. The Gaussian hill defined in \eqref{eq:hill_C} will also be used as the cost function in the numerical simulations presented later in this dissertation, \eqref{eq:hill_cost}. This provides a common cost landscape through which the analytical and numerical aspects of the problem can be considered.

\subsection{Rational Quadratic}

Considering how the Gaussian Hill describes in \eqref{eq:hill_C}, we try a simpler function for $c$ to check the Euler-ladgrange equation. We choose $c$, 
\begin{equation}
\label{eq:hill_C_rational}
c(\gamma(s)) = 
c_0+
\frac{H R^2}
{R^2+\left|\gamma(s)-p\right|^2},
\end{equation}
Where $H>0$ represents the magnitude of the increase in cost, $p\in\mathbb{R}^2$ is the centre of the hill and $R$ is the radius. The background cost is $c_0$ and the maximum cost is $c_{\max}$, then
$$
H=c_{\max}-c_0.
$$ 

The gradient of \eqref{eq:hill_C_rational} with respect to $\gamma$ is
$$
\nabla c(\gamma)
=
-\frac{2HR^2(\gamma-p)}
{\left(R^2+\left|\gamma-p\right|^2\right)^2}.
$$
Therefore, the Euler--Lagrange equation, \eqref{eq:ModelProblem_EL_Solution}, becomes
\begin{equation}
\label{eq:E_L_Rational}
\gamma^{(4)}(s)=
\frac{HR^2}{\lambda}
\frac{\gamma(s)-p}
{\left(R^2+\left|\gamma(s)-p\right|^2\right)^2}.
\end{equation}

This is a non-linear fourth-order ordinary differential equation. Although the order of the equation cannot be reduced directly from the form of the right-hand side, its non-linear behaviour can be simplified asymptotically. In particular, when
$$
\left|\gamma(s)-p\right|^2 \gg R^2,
$$
the denominator satisfies
$$
\left(R^2+\left|\gamma(s)-p\right|^2\right)^2
\sim
\left|\gamma(s)-p\right|^4.
$$
Consequently,
$$
\gamma^{(4)}(s)
\sim
\frac{HR^2}{\lambda}
\frac{\gamma(s)-p}
{\left|\gamma(s)-p\right|^4}.
$$
In the scalar case, this has the asymptotic behaviour
$$
\gamma^{(4)}(s)
\sim
\frac{C}{\left(\gamma(s)-p\right)^3} \sim \frac{1}{\gamma^3},
$$
for an appropriate constant \(C\). Thus, far from the point \(p\), the forcing term decays cubically with the distance from \(p\). This provides a much simpler ODE to solve and analyse. 

Although \eqref{eq:E_L_Rational} is a nonlinear fourth-order ordinary differential equation, its right-hand side is a smooth function of \(\gamma\). Consequently, the equation may be rewitten as a system of first-order ordinary differential equations, to which the Picard--Lindelöf theorem, Theorem 2.2~\cite{teschl2012ordinary}, applies. It follows that, for prescribed initial data
$$
\gamma(s_0),\qquad \gamma'(s_0),\qquad \gamma''(s_0),\qquad \gamma'''(s_0),
$$
there exists a unique local solution to \eqref{eq:E_L_Rational}. Thus, despite the nonlinear nature of the Euler--Lagrange equation, the corresponding initial-value problem admits a unique local solution.

{\color{blue} Will test and write a short paragraph showing it works for the numerics, and link to the jpyner notepad. might put in the numerical chapter}

\section{Functional Setting}

In this section, we investigate the properties of the integrand in \eqref{eq:ModelPrblem_s}. We begin by defining the set of admissible functions and explaining the choice of function space. We then examine the integrand to determine whether it satisfies the conditions required for coercivity, weak compactness, and weak lower semicontinuity. Finally, these properties are used within the Direct Method of the Calculus of Variations to establish the existence of a minimiser.


\subsection{Existence of minimizers}

In this section we state and prove the main result of this dissertation on existence of a minimizer of the problem we investigate. Instead of the Gaussian Hill function, we choose the cost function $c(\gamma(s))$ to be in the quadratic form. 

\subsubsection{Assumptions}




\noindent In order to make the choice of the set of the smooth curves for \eqref{eq:ModelPrblem_s} well-defined, we define the set
\begin{equation}\label{D}
\mathcal{D} = \{\gamma \in C^{2}([s_0, s_1]) \colon \gamma'' \in L^2(\mathbb{R}), c(\gamma) \in L^{1}(\mathbb{R}), |\gamma'|=1\}
\end{equation}




\subsubsection{The main result}

Before we state and prove the existence theorem for problem \eqref{eq:ModelPrblem_s} we note that it is not possible to apply the general existence theorem for Calculus of Variations given in Section 2.6.2. This is because the integrand in \eqref{eq:ModelPrblem_s} contains higher order derivative term due to the curvature. As a result the minimization problem differs fundamentally from the minimization problem in Thm. 2.44. 

There are numerous studies in the literature on minimization problems with functionals including curvature. We will refer to the ones by Boscain, Charlot and Rossi \cite{BoscainCharlotRossi}, and Bellettini \cite{Bellettini}.

\begin{theorem}
For suitably chosen boundary conditions $a, b \in \Omega$,  problem \eqref{eq:ModelPrblem_s} admits a unique minimizer in $\mathcal{D}$, which is defined in \eqref{D}, satisfying the Euler-Lagrange equation \eqref{eq:ModelProblem_EL_Solution} with $c$ chosen as in \eqref{quadratic-c}.
\end{theorem}


{\color{blue}
\textbf{I asked Chat GTP about this rational hill as a first steps/guide, is this what we'd expect to see/ want for this or should it all be ingored.}
\begin{theorem}
\label{thm:DirectMethod_RationalHill}
Let $s_0<s_1$, let $A,B\in\mathbb{R}^2$ and let $v_0,v_1\in\mathbb{R}^2$. Define the admissible set
\begin{equation}
\label{eq:admissible_rational}
\mathcal{A}=
\left\{
\gamma\in W^{2,2}((s_0,s_1);\mathbb{R}^2):
\begin{array}{l}
\gamma(s_0)=A,\quad \gamma(s_1)=B,\
\gamma'(s_0)=v_0,\quad \gamma'(s_1)=v_1
\end{array}
\right\}.
\end{equation}
Let $c_0\in\mathbb{R}$, $H>0$, $R>0$ and $\lambda>0$, and define
\begin{equation}
\label{eq:rational_hill}
c(\gamma) =  c_0+
\frac{HR^2}{R^2+|\gamma-p|^2},
\qquad p\in\mathbb{R}^2.
\end{equation}
Consider the functional
\begin{equation}
\label{eq:rational_functional}
J[\gamma]
=
\int_{s_0}^{s_1}
\left(
c(\gamma(s))
+
\lambda|\gamma''(s)|^2
+
\mu
\right),ds.
\end{equation}
Assume that $\mathcal{A}\neq\varnothing$. Then the problem
\begin{equation}
\label{eq:rational_minimisation}
\inf_{\gamma\in\mathcal{A}}J[\gamma]
\end{equation}
admits a minimiser $\overline{\gamma}\in\mathcal{A}$.
\end{theorem}

\begin{proof}
We divide the proof into three steps.

\begin{enumerate}

\item[(i)] \textit{Compactness.}

Let

$$
m=\inf_{\gamma\in\mathcal{A}}J[\gamma].
$$

Since $\mathcal{A}\neq\varnothing$, there exists $\gamma_0\in\mathcal{A}$ and hence

$$
m\leq J[\gamma_0]<\infty.
$$

Furthermore, since $H>0$ and $R>0$,

$$
c(\gamma)
=
c_0+
\frac{HR^2}{R^2+|\gamma-p|^2}
\geq c_0.
$$

Therefore,

$$
J[\gamma]
\geq
\lambda\int_{s_0}^{s_1}|\gamma''(s)|^2\,ds
+
(c_0+\mu)(s_1-s_0).
$$

Since $\lambda>0$, this provides a bound on $|\gamma''|_{L^2}$ for every minimising sequence.

Let ${\gamma_\nu}\subset\mathcal{A}$ be a minimising sequence, so that

$$
J[\gamma_\nu]\longrightarrow m
\qquad\text{as }\nu\rightarrow\infty.
$$

The preceding estimate implies that

$$
\|\gamma_\nu''\|_{L^2}
$$

is bounded. Since the endpoint values and endpoint derivatives are fixed, a second-order Poincar'e inequality gives constants $C_1,C_2>0$ such that

$$
\|\gamma_\nu\|_{W^{2,2}}
\leq
C_1\|\gamma_\nu''\|_{L^2}+C_2.
$$

Consequently, ${\gamma_\nu}$ is bounded in $W^{2,2}$.

Since $W^{2,2}((s_0,s_1);\mathbb{R}^2)$ is reflexive, there exists a subsequence, not relabelled, and some $\overline{\gamma}\in W^{2,2}((s_0,s_1);\mathbb{R}^2)$ such that

$$
\gamma_\nu\rightharpoonup\overline{\gamma}
\qquad\text{weakly in }W^{2,2}.
$$

Moreover, in one dimension the embedding

$$
W^{2,2}((s_0,s_1);\mathbb{R}^2)
\hookrightarrow C^1([s_0,s_1];\mathbb{R}^2)
$$

is compact. Hence, after passing to a further subsequence if necessary,

$$
\gamma_\nu\rightarrow\overline{\gamma}
\qquad\text{strongly in }C^1([s_0,s_1];\mathbb{R}^2).
$$

In particular, the boundary conditions are preserved, and therefore

$$
\overline{\gamma}\in\mathcal{A}.
$$

\item[(ii)] \textit{Weak lower semicontinuity.}

We show that

$$
J[\overline{\gamma}]
\leq
\liminf_{\nu\rightarrow\infty}J[\gamma_\nu].
$$

First consider the curvature term. Since

$$
\gamma_\nu''\rightharpoonup\overline{\gamma}''
\qquad\text{weakly in }L^2,
$$

and the function

$$
z\mapsto |z|^2
$$

is convex and weakly lower semicontinuous,

$$
\int_{s_0}^{s_1}|\overline{\gamma}''(s)|^2\,ds
\leq
\liminf_{\nu\rightarrow\infty}
\int_{s_0}^{s_1}|\gamma_\nu''(s)|^2\,ds.
$$

Next consider the cost term. Since

$$
\gamma_\nu\rightarrow\overline{\gamma}
\qquad\text{uniformly on }[s_0,s_1],
$$

and $c$ is continuous, we have

$$
c(\gamma_\nu(s))
\rightarrow
c(\overline{\gamma}(s))
$$

uniformly on $[s_0,s_1]$. Consequently,

$$
\int_{s_0}^{s_1}c(\gamma_\nu(s))\,ds
\longrightarrow
\int_{s_0}^{s_1}c(\overline{\gamma}(s))\,ds.
$$

Combining these two results with the fact that

$$
\int_{s_0}^{s_1}\mu\,ds
=
\mu(s_1-s_0),
$$

we obtain

$$
J[\overline{\gamma}]
\leq
\liminf_{\nu\rightarrow\infty}J[\gamma_\nu].
$$

\item[(iii)] \textit{Combine.}

Since ${\gamma_\nu}$ is a minimising sequence,

$$
J[\gamma_\nu]\longrightarrow m.
$$

By weak lower semicontinuity,

$$
J[\overline{\gamma}]
\leq
\liminf_{\nu\rightarrow\infty}J[\gamma_\nu]
=
m.
$$

However, since $\overline{\gamma}\in\mathcal{A}$ and $m$ is the infimum of $J$ over $\mathcal{A}$,

$$
J[\overline{\gamma}]\geq m.
$$

Therefore,

$$
J[\overline{\gamma}]=m.
$$

Hence $\overline{\gamma}$ is a minimiser of \eqref{eq:rational_minimisation}.
\end{enumerate}
\end{proof}
}